The declared group \(\mathcal G=S_2\times S_2\) contains the two local swap subgroups, one for each pair block. Thus the conditional construction in \texttt{def:pair-action} applies to both blocks. The action preserves blocks and arms, and each block contains exactly two members, so the four global arm-specific orbits are precisely
\[
O_{1,0},\quad O_{1,1},\quad O_{2,0},\quad O_{2,1}.
\tag{1}
\]
If \(q_{\mathrm{tar}}\) concentrates on \(v_{1,0}\), then \(T^1=(v_{1,0},1)\in O_{1,1}\) almost surely. Hence the target arm-one orbit vector assigns mass one to \(O_{1,1}\). The exclusive mate-control state \((v_{1,1},0)\) belongs to the distinct arm-zero orbit \(O_{1,0}\), so its analogue vector has mass one on \(O_{1,0}\) and zero in the target row \((1,O_{1,1})\). The corresponding equations of \texttt{def:common-selector-feasibility-lp} therefore cannot both hold. Moreover, \texttt{ass:observable-selector-kernel} excludes the control-labelled mate from a legal arm-one selector.

If a legal implemented selector instead satisfies \(q^{\mathrm{obs}}_{w,1}(O_{1,1})=0\), the only target-supported arm-one orbit has zero calibration mass. For \(k\le M\), \texttt{prop:disjoint-orbit-witness} gives
\[
\mathcal C_w=\{0\},\qquad \Pr_P(\tau_T\in\mathcal C_w\mid\mathcal T)=0.
\tag{2}
\]
For \(k=M+1\), \texttt{prop:rank-boundary} gives
\[
I_{0,w}=I_{1,w}=[-K,K],\qquad \mathcal C_w=[-2K,2K],\qquad \Pr_P(\tau_T\in\mathcal C_w\mid\mathcal T)=1.
\tag{3}
\]
This proves every claimed finite check.